% Options for packages loaded elsewhere
\PassOptionsToPackage{unicode}{hyperref}
\PassOptionsToPackage{hyphens}{url}
\documentclass[
  letterpaper,
]{article}
\usepackage{xcolor}
\usepackage[margin=0.8in]{geometry}
\usepackage{amsmath,amssymb}
\setcounter{secnumdepth}{-\maxdimen} % remove section numbering
\usepackage{iftex}
\ifPDFTeX
  \usepackage[T1]{fontenc}
  \usepackage[utf8]{inputenc}
  \usepackage{textcomp} % provide euro and other symbols
\else % if luatex or xetex
  \usepackage{unicode-math} % this also loads fontspec
  \defaultfontfeatures{Scale=MatchLowercase}
  \defaultfontfeatures[\rmfamily]{Ligatures=TeX,Scale=1}
\fi
\usepackage{lmodern}
\ifPDFTeX\else
  % xetex/luatex font selection
\fi
% Use upquote if available, for straight quotes in verbatim environments
\IfFileExists{upquote.sty}{\usepackage{upquote}}{}
\IfFileExists{microtype.sty}{% use microtype if available
  \usepackage[]{microtype}
  \UseMicrotypeSet[protrusion]{basicmath} % disable protrusion for tt fonts
}{}
\makeatletter
\@ifundefined{KOMAClassName}{% if non-KOMA class
  \IfFileExists{parskip.sty}{%
    \usepackage{parskip}
  }{% else
    \setlength{\parindent}{0pt}
    \setlength{\parskip}{6pt plus 2pt minus 1pt}}
}{% if KOMA class
  \KOMAoptions{parskip=half}}
\makeatother
\usepackage{listings}
\newcommand{\passthrough}[1]{#1}
\lstset{defaultdialect=[5.3]Lua}
\lstset{defaultdialect=[x86masm]Assembler}
\setlength{\emergencystretch}{3em} % prevent overfull lines
\providecommand{\tightlist}{%
  \setlength{\itemsep}{0pt}\setlength{\parskip}{0pt}}
% Preserve readable code, command, and table layout in Pandoc-generated review PDFs.
\usepackage{etoolbox}
\usepackage{pdflscape}
\lstset{
  basicstyle=\ttfamily\footnotesize,
  breaklines=true,
  breakatwhitespace=false,
  columns=fullflexible,
  keepspaces=true,
  showstringspaces=false,
  upquote=true
}
\AtBeginEnvironment{longtable}{\footnotesize}
\renewcommand{\arraystretch}{1.08}
\setlength{\emergencystretch}{3em}
\usepackage{bookmark}
\IfFileExists{xurl.sty}{\usepackage{xurl}}{} % add URL line breaks if available
\urlstyle{same}
% fallback for those not using the hyperref driver hyperxmp:
\makeatletter
\@ifundefined{xmpquote}{\newcommand{\xmpquote}[1]{#1}}{}
\makeatother
\hypersetup{
  pdftitle={Password Bonobo Repository Foundation and Compatibility Dossier Specification},
  hidelinks,
  pdfcreator={LaTeX via pandoc}}

\title{Password Bonobo Repository Foundation and Compatibility Dossier Specification}
\author{}
\date{}

\begin{document}
\maketitle

{
\setcounter{tocdepth}{3}
\tableofcontents
}
\section{Password Bonobo Repository Foundation and Compatibility Dossier
Specification}\label{password-bonobo-repository-foundation-and-compatibility-dossier-specification}

Date: 2026-08-22

Status: Approved program design; implementation plan pending execution

Parent design: \href{./password-bonobo-python-reimplementation-design.md}{Password Bonobo Python Reimplementation}

\subsection{1. Purpose}\label{1-purpose}

This subproject establishes the Password Bonobo repository, its enforceable engineering policy, and a neutral
compatibility record for the selected Password Gorilla baseline. It creates the boundary that later subprojects will use
to implement original Python code without consulting or translating Gorilla source during product development.

This subproject produces infrastructure, documentation, and an empty typed Python package boundary. It does not
implement PasswordSafe parsing, vault behavior, URL auditing, synchronization, or a user interface.

\subsection{2. Required Outcomes}\label{2-required-outcomes}

The subproject is complete when it has:

\begin{enumerate}
\def\labelenumi{\arabic{enumi}.}
\tightlist
\item
  Established a Git repository with an intentional tracked-file boundary and a reviewable commit history.
\item
  Defined the GPL-3.0-or-later project license and documented the unresolved iOS distribution-exception work.
\item
  Configured Python 3.14, strict typing, formatting, linting, tests, security checks, license checks, and continuous
  integration.
\item
  Enforced the project-specific Python source documentation and spacing rules.
\item
  Preserved a selected Gorilla revision in an untouched checkout outside the Bonobo repository.
\item
  Recorded the upstream URL, exact commit, retrieval instructions, observed license, and source-analysis boundary.
\item
  Produced a neutral behavior dossier, feature-parity matrix, and executable test-oracle catalog.
\item
  Created durable project memory that records decisions, status, verification evidence, and the next approved work.
\item
  Demonstrated that generated documentation is reproducible and that generated PDFs remain uncommitted.
\end{enumerate}

\subsection{3. Repository Boundary}\label{3-repository-boundary}

The repository root is the current \passthrough{\lstinline!Password Bonobo!} workspace. The initial tracked set will
include:

\begin{itemize}
\tightlist
\item
  Root policy and contributor documents.
\item
  Repository configuration and quality-tool configuration.
\item
  Bonobo-authored specifications, plans, compatibility documents, and project memory.
\item
  The empty \passthrough{\lstinline!bonobo\_core!} package boundary and its foundation tests.
\item
  Cross-platform validation tools owned by Bonobo.
\item
  Continuous-integration workflows.
\end{itemize}

The initial tracked set will exclude:

\begin{itemize}
\tightlist
\item
  The locally mirrored shared standards under \passthrough{\lstinline!docs/prompts/!}.
\item
  Generated PDFs and temporary render files.
\item
  Logs, caches, local environments, build outputs, and editor state.
\item
  PasswordSafe databases except explicitly allowlisted synthetic fixtures added by a later subproject.
\item
  The Gorilla checkout and every other upstream source tree.
\item
  Real credentials, tokens, personal vaults, or provider metadata.
\end{itemize}

The repository will use \passthrough{\lstinline!main!} as its primary branch. Foundation work will occur on
\nolinkurl{foundation/compatibility-dossier} after the initial baseline commit.

\subsection{4. Upstream Gorilla Baseline}\label{4-upstream-gorilla-baseline}

The selected compatibility reference is:

\begin{itemize}
\tightlist
\item
  Repository: \nolinkurl{https://github.com/zdia/gorilla.git}
\item
  Branch observed: \passthrough{\lstinline!master!}
\item
  Commit: \nolinkurl{6728e85c05ac25357b8f19f541487b9d26a97402}
\item
  Observation date: 2026-08-22
\item
  Declared license observed in the upstream README: GPL-2.0-or-later
\end{itemize}

The default research location will be the sibling directory
\passthrough{\lstinline!../Password Bonobo Research/gorilla!}. An operator may choose another location outside the
Bonobo repository, but the recorded commit and verification procedure do not change.

The checkout will be detached at the selected commit and treated as read-only. Updating the reference requires a
separate review, a new pinned commit, a dossier delta, and an explicit provenance entry.

\subsection{5. Source-Analysis Boundary}\label{5-source-analysis-boundary}

Gorilla source is evidence about behavior, not implementation material for Bonobo.

The research pass may inspect:

\begin{itemize}
\tightlist
\item
  \passthrough{\lstinline!sources/gorilla.tcl!} for application workflows and user-visible behavior.
\item
  \passthrough{\lstinline!sources/pwsafe/*.tcl!} for supported data concepts and compatibility behavior.
\item
  \passthrough{\lstinline!sources/help.txt!}, message catalogs, and upstream README material for user-facing contracts.
\item
  \passthrough{\lstinline!unit-tests/**/*.test!} and \passthrough{\lstinline!unit-tests/**/*.tcl!} for behavioral
  examples and edge cases.
\item
  Upstream change history for compatibility-relevant corrections.
\end{itemize}

The research pass must not:

\begin{itemize}
\tightlist
\item
  Modify the Gorilla checkout.
\item
  Copy source, comments, identifiers, file organization, control flow, UI assets, or translations into Bonobo.
\item
  Import upstream \passthrough{\lstinline!.psafe3!} files as Bonobo fixtures. Upstream vault files remain external
  research-only and may never enter tracked fixture paths; every tracked credential or vault fixture requires recorded
  synthetic-data proof, license review, sensitive-data scanning, and approval.
\item
  Add Python type hints or Bonobo comments to Tcl files.
\item
  Cite source fragments in product code.
\item
  Treat the absence of an upstream test as proof that a behavior is unsupported.
\end{itemize}

The dossier will express behavior in neutral domain language. It may cite an upstream path and line range as evidence,
but it will not reproduce source text. Later product implementation will work from the approved Bonobo specifications,
the PasswordSafe format specification, the dossier, and synthetic tests.

\subsection{6. Compatibility Deliverables}\label{6-compatibility-deliverables}

\subsubsection{6.1 Upstream baseline record}\label{61-upstream-baseline-record}

\nolinkurl{docs/compatibility/gorilla/upstream-baseline.md} will record the selected revision, retrieval and
verification commands, license observation, external checkout rule, evidence conventions, and update procedure.

\subsubsection{6.2 Behavior dossier}\label{62-behavior-dossier}

\nolinkurl{docs/compatibility/gorilla/behavior-dossier.md} will describe:

\begin{itemize}
\tightlist
\item
  Vault creation, opening, authentication, saving, closing, locking, and recovery behavior.
\item
  Entry creation, editing, deletion, protection, history, aliases, and shortcuts.
\item
  Groups, tree operations, search, filtering, and selection behavior.
\item
  Password generation and policy behavior.
\item
  Clipboard, browser launch, and AutoType behavior.
\item
  Import, export, merge, backup, preference, and recent-file behavior.
\item
  PasswordSafe version handling and user-authored metadata behavior.
\item
  Errors, confirmations, destructive-action protections, and observable edge cases.
\item
  Platform-specific behavior and known limitations.
\end{itemize}

Each behavioral statement will carry an evidence identifier. Evidence records will include the upstream revision,
relative source or test path, line range or test name, evidence kind, and a neutral observation. No copied code will be
included.

\subsubsection{6.3 Feature-parity matrix}\label{63-feature-parity-matrix}

\nolinkurl{docs/compatibility/gorilla/feature-parity-matrix.md} will use these statuses:

\begin{itemize}
\tightlist
\item
  \passthrough{\lstinline!Required!}: meaningful Gorilla compatibility required for stable Bonobo.
\item
  \passthrough{\lstinline!Modernized!}: the behavior is required but its interaction design will change.
\item
  \passthrough{\lstinline!Deferred!}: accepted for a named later Bonobo subproject.
\item
  \passthrough{\lstinline!Excluded!}: intentionally unsupported with an approved rationale.
\item
  \passthrough{\lstinline!Bonobo extension!}: behavior added by Bonobo rather than inherited from Gorilla.
\item
  \passthrough{\lstinline!Unverified!}: evidence is insufficient and additional research is required.
\end{itemize}

Every row will identify its dossier evidence, Bonobo owner subproject, intended platforms, data-loss relevance, security
relevance, and acceptance-test identifier.

\subsubsection{6.4 Test-oracle catalog}\label{64-test-oracle-catalog}

\nolinkurl{docs/compatibility/gorilla/test-oracles.md} will define black-box scenarios using synthetic data. Each
scenario will state setup, action, observable result, cleanup, supported evidence, and whether the result must be
confirmed against Gorilla, Password Safe, or both.

The catalog will distinguish:

\begin{itemize}
\tightlist
\item
  Normative PasswordSafe requirements.
\item
  Observed Gorilla compatibility behavior.
\item
  Bonobo product decisions.
\item
  Unresolved questions that block a compatibility claim.
\end{itemize}

\subsection{7. Python Foundation}\label{7-python-foundation}

The repository will target CPython \passthrough{\lstinline!>=3.14,<3.15!} for the first development baseline. A
version-floor change requires platform and dependency review because the same core is intended for desktop, Android, and
iOS.

The initial package boundary will be \passthrough{\lstinline!src/bonobo\_core/!}. It will contain no vault
implementation in this subproject. The package will be marked typed and expose only version metadata.

The project will use:

\begin{itemize}
\tightlist
\item
  Hatchling as the Python build backend.
\item
  \passthrough{\lstinline!uv!} for Python installation, dependency locking, and reproducible command execution.
\item
  autopep8 as the formatter, configured to preserve required three-blank-line spacing.
\item
  Ruff as the linter and import-order checker.
\item
  mypy in strict mode as the required static type checker.
\item
  pytest and coverage for tests.
\item
  Bandit and pip-audit for source and dependency security checks.
\item
  REUSE tooling for machine-checkable license metadata.
\end{itemize}

A repository test will enforce the required module docstrings, \passthrough{\lstinline!\#\#\#\#!} declaration blocks,
decorator placement, and three-blank-line separation for maintained Python files. Tool configuration must not silently
override those rules.

\subsection{8. Continuous Integration}\label{8-continuous-integration}

The foundation workflow will run on Windows, macOS, and Linux with Python 3.14. Each job will:

\begin{enumerate}
\def\labelenumi{\arabic{enumi}.}
\tightlist
\item
  Check out the Bonobo repository without any Gorilla source.
\item
  Install the pinned Python toolchain from the lock file.
\item
  Check formatting.
\item
  Run Ruff and strict mypy.
\item
  Run pytest with coverage.
\item
  Run the Python structure-policy test.
\item
  Run Bandit, pip-audit, and REUSE checks.
\item
  Verify that forbidden file types and likely credential artifacts are absent from tracked files.
\item
  Build the Python wheel and source distribution.
\end{enumerate}

The workflow must not clone Gorilla, contact password providers, or require secrets.

\subsection{9. Licensing and Contributions}\label{9-licensing-and-contributions}

Bonobo-authored code will use SPDX identifier \passthrough{\lstinline!GPL-3.0-or-later!}. The repository will contain
the complete corresponding license text and machine-readable REUSE metadata.

The App Store distribution exception is not finalized by this subproject. The repository will document its purpose,
scope constraints, dependency implications, and required legal review. It will not publish provisional exception text as
if it were approved.

External contributions will remain disabled until the contribution terms can preserve both GPL-3.0-or-later licensing
and the ability to publish an approved distribution exception for Bonobo-authored iOS code.

\subsection{10. Project Memory}\label{10-project-memory}

Durable memory will record:

\begin{itemize}
\tightlist
\item
  The approved program and subproject specifications.
\item
  The pinned Gorilla revision and source boundary.
\item
  Toolchain and repository decisions with rationale.
\item
  Current subproject status and verification evidence.
\item
  Known risks and unresolved decisions.
\item
  The exact next approved subproject.
\end{itemize}

Memory will not contain credentials, sensitive test data, copied Gorilla source, or transient command transcripts.

\subsection{11. Verification}\label{11-verification}

Foundation verification will include:

\begin{itemize}
\tightlist
\item
  A clean \passthrough{\lstinline!git status!} after each committed task, excluding intentionally untracked generated
  PDFs.
\item
  A test proving the Python package is importable and typed.
\item
  Passing formatter, Ruff, strict mypy, pytest, coverage, Bandit, pip-audit, and REUSE checks.
\item
  A negative structure-policy fixture proving undocumented declarations are rejected.
\item
  A tracked-file audit proving the Gorilla tree, PasswordSafe databases, PDFs, logs, and shared standards are excluded.
\item
  A detached upstream checkout whose \passthrough{\lstinline!HEAD!} equals the selected commit and whose worktree is
  clean.
\item
  Dossier evidence coverage for every initial parity-matrix row.
\item
  A scan proving the dossier contains no copied multi-line source fragments or unresolved drafting tokens.
\item
  Successful Markdown-to-LaTeX generation and XeLaTeX compilation for substantial project documents.
\item
  Page-by-page visual inspection of generated PDFs.
\end{itemize}

\subsection{12. Acceptance Criteria}\label{12-acceptance-criteria}

This subproject is accepted when:

\begin{enumerate}
\def\labelenumi{\arabic{enumi}.}
\tightlist
\item
  The repository boundary is explicit, enforced, and free of imported Gorilla implementation material.
\item
  The selected Gorilla revision can be independently retrieved and verified from the baseline record.
\item
  The upstream checkout remains untouched outside the Bonobo repository.
\item
  The behavior dossier and parity matrix cover every meaningful feature family found during the research pass.
\item
  Every parity row links to evidence and an acceptance-test identifier or is explicitly
  \passthrough{\lstinline!Unverified!}.
\item
  Later implementers can work from Bonobo documents without reading Gorilla source.
\item
  The Python foundation passes all required local and continuous-integration gates.
\item
  Licensing and provenance checks pass without implying that the iOS exception is already approved.
\item
  Project memory identifies the lossless PasswordSafe core as the next subproject.
\item
  Markdown and LaTeX sources are committed while generated PDFs remain uncommitted.
\end{enumerate}

\subsection{13. Known Risks}\label{13-known-risks}

\begin{itemize}
\tightlist
\item
  A single researcher can unintentionally preserve upstream expression in notes; neutral prose and evidence-only
  citations reduce but do not eliminate that risk.
\item
  Some Gorilla behavior may depend on Tcl/Tk, operating-system state, or unavailable legacy packages and may initially
  remain \passthrough{\lstinline!Unverified!}.
\item
  Upstream test databases may contain data unsuitable for redistribution and therefore cannot be adopted by default.
\item
  Python 3.14 support can vary among later binary dependencies and must be rechecked before each platform subproject.
\item
  App Store rules and dependency licenses can change before the iOS release.
\end{itemize}

These risks do not relax the no-copy boundary or the no-loss compatibility requirement.

\subsection{14. References}\label{14-references}

\begin{itemize}
\tightlist
\item
  Password Bonobo program design: \nolinkurl{docs/specs/password-bonobo-python-reimplementation-design.md}
\item
  Password Bonobo URL-audit design: \nolinkurl{docs/specs/password-bonobo-url-audit-design.md}
\item
  Password Gorilla: \url{https://github.com/zdia/gorilla}
\item
  PasswordSafe V3 format: \url{https://github.com/pwsafe/pwsafe/blob/master/docs/formatV3.txt}
\item
  Python on Android: \url{https://docs.python.org/3.14/using/android.html}
\item
  Python on iOS: \url{https://docs.python.org/3.14/using/ios.html}
\item
  REUSE specification: \url{https://reuse.software/spec/}
\end{itemize}

\end{document}
